\chapter{Origin, Scope, and Acknowledgement}
\label{ch:origin}

\section{Initial Concept}

The initial concept selected a broad read-only capability surface combining
geospatial and statistical queries. This was treated as a practical stress test
for MCP delivery in public-sector contexts \cite{apps_to_answers_brief,deep_research_report}.

\section{Acknowledgement of Early Contribution}

The repository history includes an explicit linkage to early \texttt{os-mcp}
work. The origin narrative also records acknowledgement of early contribution by
Chris Carlon in initial project prompts captured by Codex reporting artifacts
\cite{codex_long_horizon_20260304}.

\section{Change of Course}

Delivery moved toward a fully tracked repository model with:
\begin{itemize}
  \item durable context tracking,
  \item progress tracker workstreams,
  \item release-note discipline,
  \item test and troubleshooting evidence loops.
\end{itemize}

\begin{figure}[h]
\centering
\begin{tikzpicture}[
  node distance=1.7cm,
  box/.style={draw, rounded corners, align=center, inner sep=6pt},
  >=Latex
]
\node[box] (a) {Initial AI\\design stage};
\node[box, right=of a] (b) {Early \texttt{os-mcp}\\phase};
\node[box, right=of b] (c) {New tracked\\repository path};
\node[box, right=of c] (d) {\texttt{mcp-geo}\\delivery stream};
\draw[->] (a) -- (b);
\draw[->] (b) -- (c);
\draw[->] (c) -- (d);
\end{tikzpicture}
\caption{Observed project lineage represented in repository artifacts.}
\end{figure}
